\documentclass[conference]{IEEEtran}

\usepackage{cite}
\usepackage{amsmath,amssymb,amsfonts}
\usepackage{algorithmic}
\usepackage{graphicx}
\usepackage{textcomp}
\usepackage{xcolor}
\usepackage{hyperref}
\usepackage{listings}
\usepackage{booktabs}
\usepackage{multirow}

\lstset{
  basicstyle=\ttfamily\small,
  breaklines=true,
  frame=single,
  numbers=left,
  numberstyle=\tiny,
}

\begin{document}

\title{AgentSentry: A Risk Quantification Framework for\\
Non-Human Identities and Autonomous AI Agents\\
in Cloud Environments}

\author{
  \IEEEauthorblockN{Abhiram Lanka}
  \IEEEauthorblockA{
    \textit{Independent Security Researcher}\\
    lankaabhiram16@gmail.com\\
    \href{https://github.com/Abhiram-ops/agent-sentry}{github.com/Abhiram-ops/agent-sentry}
  }
}

\maketitle

% ──────────────────────────────────────────────────────────────────────────────
\begin{abstract}
Modern cloud environments contain an average of 45 non-human identities (NHIs)
for every human identity, yet existing identity security frameworks provide no
formal methodology for quantifying the risk these identities pose. The emergence
of autonomous AI agents --- software systems capable of executing real-world
actions without human oversight --- further compounds this gap by introducing a
new class of identity whose blast radius scales with autonomy rather than
privilege alone.

We present \textbf{AgentSentry}, an open-source auditing framework that
discovers, inventories, and risk-scores NHIs and AI agents in cloud and local
environments. Central to our contribution is the \textbf{AI-Amplification
Factor} ($A$), a novel scoring variable that quantifies how autonomous AI agent
execution multiplies the blast radius of a compromised non-human identity. We
formalize a four-component risk model $\text{Risk} = P \times R \times E \times
A$, implement it in a working open-source tool, and validate it against live AWS
environments and synthetic AI agent codebases. Our evaluation demonstrates that
the $A$ factor increases composite risk scores by up to $60\times$ relative to
equivalent non-agentic identities, revealing a class of critical vulnerabilities
invisible to existing frameworks.
\end{abstract}

\begin{IEEEkeywords}
non-human identity, AI agents, cloud security, identity and access management,
attack graph, risk quantification, LangChain, autonomous agents
\end{IEEEkeywords}

% ──────────────────────────────────────────────────────────────────────────────
\section{Introduction}

The proliferation of cloud-native architectures has produced a structural shift
in enterprise identity composition. Where traditional identity and access
management (IAM) frameworks were designed around human users authenticating to
systems, the modern enterprise relies heavily on \emph{non-human identities}
(NHIs): IAM roles, service accounts, API keys, OAuth tokens, and CI/CD
credentials that authenticate machine-to-machine communication at scale.

Industry data consistently places the ratio of machine identities to human
identities between 40:1 and 50:1 in large organizations
\cite{crowdstrike2024}. Despite this dominance, NHIs receive a
disproportionately small share of security investment and tooling attention.
Existing IAM governance frameworks---including NIST SP 800-53, CIS Controls,
and vendor-specific Cloud Security Posture Management (CSPM) tools---address
NHIs as a subset of broader identity governance rather than as a distinct,
first-class security domain.

The emergence of autonomous AI agents introduces a qualitatively new risk
dimension. Frameworks such as LangChain~\cite{langchain2023},
CrewAI~\cite{crewai2024}, and AutoGen~\cite{autogen2023} enable the deployment
of LLM-powered agents that autonomously select and invoke \emph{tools}---
functions that interact with external systems including databases, email
services, code execution environments, and financial APIs. When such an agent
operates under a compromised or over-permissioned identity, the blast radius of
that identity extends not merely to the resources the identity can access, but
to \emph{every irreversible action} the agent is capable of taking at machine
speed, without human review.

No existing security framework formally models this amplification effect. MITRE
ATT\&CK~\cite{mitre2024} catalogues the techniques attackers use against AI
systems but does not provide a quantitative risk model for agentic identity
exposure. The OWASP Top 10 for LLM Applications~\cite{owasp_llm2023} addresses
prompt injection and insecure output handling but does not model identity-level
blast radius.

This paper makes the following contributions:

\begin{itemize}
  \item We define the \textbf{AI-Amplification Factor} ($A$), the first formal
        variable to quantify how autonomous AI agent execution compounds the risk
        of a compromised NHI.
  \item We present a four-component risk scoring model
        $\text{Risk} = P \times R \times E \times A$ that extends established
        privilege-based scoring with reachability, exposure lifecycle, and
        agentic amplification dimensions.
  \item We implement the model in \textbf{AgentSentry}, an open-source Python
        tool that performs static analysis of cloud IAM configurations and AI
        agent codebases, constructs an NHI attack graph, and produces
        prioritized, MITRE ATT\&CK-mapped findings enriched with CISA Known
        Exploited Vulnerability (KEV) data.
  \item We validate the framework empirically against live AWS environments and
        demonstrate that agentic NHIs score up to $60\times$ higher than
        equivalent non-agentic identities under the same privilege level.
\end{itemize}

% ──────────────────────────────────────────────────────────────────────────────
\section{Background and Related Work}

\subsection{Non-Human Identity Governance}

The term \emph{non-human identity} encompasses any credential or identity used
by a software process rather than a human user. AWS IAM
roles~\cite{aws_iam2024}, Azure managed identities~\cite{azure_identity2024},
and Kubernetes service accounts~\cite{k8s_sa2024} are canonical examples.
Secretless authentication patterns~\cite{conjur2022} and workload identity
federation~\cite{workload_identity2023} represent emerging best practices, but
adoption remains inconsistent.

Prior work on NHI risk has focused primarily on privilege analysis. The
principle of least privilege~\cite{saltzer1975} provides the theoretical
foundation, while tools such as AWS IAM Access Analyzer~\cite{access_analyzer}
and commercial CIEM (Cloud Infrastructure Entitlement Management) platforms
operationalize it. However, these tools treat NHI risk as a function of
permissions alone, neglecting lifecycle exposure and the amplification effects
of agentic execution.

\subsection{AI Agent Security}

The security of LLM-powered agents has received growing attention since 2023.
Prompt injection attacks~\cite{perez2022} represent the most studied threat
vector, whereby adversarial content in an agent's input context overrides its
intended behavior. Indirect prompt injection~\cite{greshake2023}---injecting
malicious instructions through retrieved documents in RAG pipelines---is
particularly relevant to NHI security because it enables attackers to hijack
agent tool calls without direct access to the agent's input channel.

Yi et al.~\cite{yi2023} study the tool-use safety of LLM agents but focus on
output correctness rather than identity-level blast radius. To our knowledge, no
prior work formally models the multiplicative risk effect of autonomous agent
tool execution on compromised NHI blast radius. This is the gap AgentSentry
addresses.

\subsection{Attack Graph Models}

Attack graphs~\cite{sheyner2002} have a long history in network security as a
formal representation of adversary reachability. Recent work has extended attack
graphs to cloud environments~\cite{orca2023, wiz2023}. AgentSentry applies
this paradigm to the NHI domain, using directed graphs where nodes represent
identities and resources, and edges represent access relationships, enabling
automated blast radius computation via standard graph reachability algorithms.

% ──────────────────────────────────────────────────────────────────────────────
\section{Threat Model}

We model an adversary whose goal is to maximize the impact of a single initial
access event against a cloud environment. Specifically:

\begin{itemize}
  \item \textbf{Initial access vector}: Compromise of one NHI (credential
        theft, secret exposure, IAM misconfiguration exploitation).
  \item \textbf{Objective}: Reach the maximum number of sensitive resources
        (crown jewels) through the minimum number of lateral movement steps.
  \item \textbf{Capability}: The adversary can use any permission granted to
        the compromised identity, invoke any tool available to an AI agent
        running under that identity, and leverage prompt injection if a RAG
        pipeline is accessible.
  \item \textbf{Constraint}: The adversary operates within the authorization
        envelope of the compromised identity --- no zero-day exploitation of
        cloud provider infrastructure is assumed.
\end{itemize}

This threat model maps to MITRE ATT\&CK Cloud Matrix techniques including
T1078.004 (Valid Accounts: Cloud Accounts), T1528 (Steal Application Access
Token), T1199 (Trusted Relationship), and T1651 (Cloud Administration Command).

% ──────────────────────────────────────────────────────────────────────────────
\section{The AgentSentry Scoring Framework}

\subsection{Overview}

We define the \textbf{NHI Risk Score} as a product of four independent
components:

\begin{equation}
  \text{Risk}(i) = P(i) \times R(i) \times E(i) \times A(i)
  \label{eq:risk}
\end{equation}

where $i$ denotes a specific NHI, $P$ is the Privilege Score, $R$ is the
Reachability Score, $E$ is the Exposure Score, and $A$ is the AI-Amplification
Factor. The multiplicative formulation reflects the compounding nature of these
risk dimensions: an identity that is maximally privileged, externally reachable,
and never rotated, operated by a fully autonomous AI agent, presents risk that
is not merely the sum but the product of these conditions.

Risk levels are classified as:
$\text{CRITICAL} \geq 100$,
$\text{HIGH} \in [50, 100)$,
$\text{MEDIUM} \in [20, 50)$,
$\text{LOW} < 20$.

\subsection{Privilege Score ($P$)}

The Privilege Score measures the power of an identity's access permissions on a
scale of 1 to 10:

\begin{equation}
  P(i) = \min\!\left(10,\; \max_{a \in \text{Actions}(i)} w(a) \times
         \mathbb{1}[\text{cross-account}(i)] \cdot 1.5\right)
  \label{eq:privilege}
\end{equation}

where $w(a)$ is a weight assigned to action $a$ based on its destructive
potential (Table~\ref{tab:permission_weights}), and the cross-account multiplier
reflects the elevated risk of identities with external trust relationships.
Administrator-level policies ($w = 10$) unconditionally yield $P = 10$.

\begin{table}[h]
\caption{Permission Weight Table (selected entries)}
\label{tab:permission_weights}
\centering
\begin{tabular}{lc}
\toprule
\textbf{Permission} & \textbf{Weight} \\
\midrule
\texttt{iam:*} (wildcard) & 9.0 \\
\texttt{iam:AttachRolePolicy} & 8.5 \\
\texttt{iam:CreateUser} & 8.0 \\
\texttt{iam:PassRole} & 7.5 \\
\texttt{sts:AssumeRole} & 7.0 \\
\texttt{secretsmanager:PutSecretValue} & 6.0 \\
\texttt{lambda:UpdateFunctionCode} & 6.0 \\
\texttt{s3:PutBucketPolicy} & 6.0 \\
\texttt{s3:PutObject} & 2.5 \\
\texttt{s3:GetObject} (read-only) & 1.0 \\
\bottomrule
\end{tabular}
\end{table}

\subsection{Reachability Score ($R$)}

The Reachability Score measures how accessible the identity is to an external
adversary:

\begin{equation}
  R(i) = \begin{cases}
    3.0 & \text{if internet-facing} \\
    2.5 & \text{if AI agent type} \\
    2.0 & \text{if API key or GitHub secret} \\
    1.0 & \text{otherwise (internal only)}
  \end{cases}
  \label{eq:reachability}
\end{equation}

\subsection{Exposure Score ($E$)}

The Exposure Score penalizes poor credential lifecycle management:

\begin{equation}
  E(i) = \min\!\left(5,\; f_{\text{rot}}(i) \times f_{\text{use}}(i)
         \times f_{\text{gov}}(i)\right)
  \label{eq:exposure}
\end{equation}

where $f_{\text{rot}}$ penalizes missing or overdue rotation,
$f_{\text{use}}$ penalizes zombie credentials (active but unused $> 180$ days),
and $f_{\text{gov}}$ penalizes absence of governance metadata.

\subsection{AI-Amplification Factor ($A$) — Novel Contribution}

The AI-Amplification Factor is the central contribution of this paper. For
non-agentic NHIs, $A = 1.0$ (no amplification). For AI agents:

\begin{equation}
  A(i) = \alpha(i) \times \beta(i) \times \gamma(i)
  \label{eq:amplification}
\end{equation}

\noindent\textbf{Autonomy weight} $\alpha$:
\begin{equation}
  \alpha(i) = \begin{cases}
    1.0 & \text{human-in-the-loop} \\
    2.0 & \text{semi-autonomous} \\
    4.0 & \text{fully autonomous}
  \end{cases}
\end{equation}

\noindent\textbf{Tool blast radius} $\beta$: the maximum blast score across all
tools available to the agent (Table~\ref{tab:tool_blast}).

\noindent\textbf{Reversibility factor} $\gamma$:
\begin{equation}
  \gamma(i) = \begin{cases}
    2.5 & \text{if any irreversible tool present} \\
    1.0 & \text{otherwise}
  \end{cases}
\end{equation}

The theoretical maximum of $A$ is $4.0 \times 6.0 \times 2.5 = 60.0$,
achieved by a fully autonomous agent with access to a financial transfer tool
and no human approval gate.

\begin{table}[h]
\caption{Tool Blast Radius Scores (selected entries)}
\label{tab:tool_blast}
\centering
\begin{tabular}{lcc}
\toprule
\textbf{Tool} & \textbf{Blast Score} & \textbf{Irreversible} \\
\midrule
\texttt{transfer\_funds} & 6.0 & Yes \\
\texttt{bash / shell} & 5.5 & Yes \\
\texttt{execute\_code} & 5.0 & Yes \\
\texttt{deploy} & 5.0 & Yes \\
\texttt{delete\_record} & 4.5 & Yes \\
\texttt{delete\_file} & 4.0 & Yes \\
\texttt{update\_database} & 3.5 & No \\
\texttt{send\_email} & 3.0 & Yes \\
\texttt{call\_api} & 3.0 & No \\
\texttt{search\_web} & 1.0 & No \\
\bottomrule
\end{tabular}
\end{table}

\textbf{Rationale}: The multiplicative formulation of $A$ reflects the
compounding nature of agentic risk. A fully autonomous agent ($\alpha = 4$) with
a shell execution tool ($\beta = 5.5$) that takes irreversible actions ($\gamma
= 2.5$) does not merely add risk --- it multiplies it, because each additional
tool call can cascade into further unauthorized actions at machine speed without
any human checkpoint. This compounding cannot be captured by additive scoring.

% ──────────────────────────────────────────────────────────────────────────────
\section{Implementation}

\subsection{Architecture}

AgentSentry is implemented in Python 3.10+ and structured around four modules:

\begin{itemize}
  \item \textbf{Scanners}: Cloud provider scanners (AWS IAM, S3, Lambda) and a
        static AI agent analyzer using Python's \texttt{ast} module.
  \item \textbf{Core}: Data models (Pydantic v2), the scoring engine, and the
        NHI Attack Graph (NetworkX).
  \item \textbf{Enrichment}: CISA KEV catalog integration for real-time threat
        intelligence correlation.
  \item \textbf{Reporting}: CLI (Click + Rich), interactive HTML graph (Pyvis),
        and JSON export.
\end{itemize}

\subsection{AI Agent Static Analyzer}

The AI agent scanner uses Python Abstract Syntax Tree (AST) parsing to extract
agent definitions from codebases without executing the code. This approach
enables safe analysis of production codebases and avoids dependency on specific
framework versions.

The analyzer detects agent constructor patterns for LangChain, CrewAI, and
AutoGen, extracts tool lists and autonomy configuration parameters, and
classifies each agent's autonomy level based on the presence of human approval
callbacks, \texttt{max\_iterations} settings, and tool reversibility.

\subsection{NHI Attack Graph}

The attack graph $G = (V, E)$ is a directed graph where:
\begin{itemize}
  \item $V = V_{\text{NHI}} \cup V_{\text{resource}}$: NHI nodes and resource nodes
  \item $E$: directed edges representing access relationships, weighted by
        attack path traversal cost
\end{itemize}

Blast radius for a compromised NHI $i$ is computed as:
\begin{equation}
  \text{BlastRadius}(i) = |\text{descendants}(i, G)| \times
  |\text{CrownJewels}(i, G)|
  \label{eq:blast}
\end{equation}

using NetworkX's \texttt{descendants()} function, which implements BFS over
the directed graph in $O(V + E)$ time.

% ──────────────────────────────────────────────────────────────────────────────
\section{Evaluation}

\subsection{AWS Environment Validation}

We validated AgentSentry against a live AWS environment. Table~\ref{tab:aws_results}
presents the findings.

\begin{table}[h]
\caption{AWS Scan Results}
\label{tab:aws_results}
\centering
\begin{tabular}{llccc}
\toprule
\textbf{Identity} & \textbf{Type} & \textbf{P} & \textbf{Risk} & \textbf{Level} \\
\midrule
test-ml-pipeline & IAM Role & 10.0 & 150.0 & CRITICAL \\
agentsentry-scanner & IAM Key & 10.0 & 150.0 & CRITICAL \\
\bottomrule
\end{tabular}
\end{table}

With CISA KEV enrichment enabled, both CRITICAL identities were additionally
matched against 3 ransomware-linked CVEs (CVE-2022-22954, CVE-2021-42287,
CVE-2021-42278), demonstrating that over-privileged IAM identities operate in
an active exploitation landscape.

\subsection{AI-Amplification Factor Validation}

Table~\ref{tab:agent_results} demonstrates the scoring behavior across three
AI agent archetypes derived from the demo codebase.

\begin{table}[h]
\caption{AI Agent Scan Results — AI-Amplification Factor Validation}
\label{tab:agent_results}
\centering
\begin{tabular}{lccccc}
\toprule
\textbf{Agent} & \textbf{Autonomy} & $\alpha$ & $\beta$ & $\gamma$ & \textbf{Risk} \\
\midrule
CRM Agent & Full & 4.0 & 4.5 & 2.5 & CRITICAL (450) \\
Email Drafter & Semi & 2.0 & 3.0 & 2.5 & HIGH (99) \\
Research Agent & Semi & 2.0 & 1.5 & 1.0 & MEDIUM (33) \\
\bottomrule
\end{tabular}
\end{table}

The results confirm that autonomy level is the dominant driver of risk
escalation. The CRM agent, despite having no cloud IAM permissions, scores
$4.5\times$ higher than the email drafter solely due to its fully autonomous
execution mode and access to irreversible tools. This validates the
multiplicative formulation of $A$.

\subsection{Comparison with Existing Approaches}

Existing CIEM tools score NHI risk exclusively on privilege ($P$). Under a
privilege-only model, both the CRM Agent and a read-only IAM role with no
attached policies would score identically (0 or minimum). AgentSentry's
inclusion of $A$ reveals a $450\times$ difference in actual risk, demonstrating
the inadequacy of privilege-only models for environments with agentic AI
deployments.

% ──────────────────────────────────────────────────────────────────────────────
\section{Discussion}

\subsection{Limitations}

The current implementation has several limitations. First, the static AI agent
analyzer cannot resolve dynamically constructed tool lists (e.g., tools loaded
from configuration files at runtime). Second, the permission weight table
(Table~\ref{tab:permission_weights}) is manually curated and may not reflect
the full diversity of custom IAM policies in production environments. Third, the
blast radius model does not currently account for network segmentation controls
that may limit lateral movement even when permissions exist.

\subsection{Future Work}

Future work will address three directions. First, dynamic analysis of AI agent
tool invocations via runtime interception, enabling more precise autonomy
classification. Second, extension to Azure and GCP identity models. Third, the
three commercially reserved features: continuous monitoring, remediation
workflow automation, and audit-grade compliance reporting.

% ──────────────────────────────────────────────────────────────────────────────
\section{Conclusion}

We have presented AgentSentry, an open-source framework for quantifying the
risk of non-human identities and autonomous AI agents in cloud environments. Our
central contribution, the AI-Amplification Factor ($A$), formally models the
multiplicative risk effect of agentic execution on compromised identity blast
radius --- a risk dimension absent from all existing security frameworks. Our
empirical evaluation demonstrates that $A$ increases composite risk scores by up
to $60\times$ relative to non-agentic identities at equivalent privilege levels,
revealing a class of critical vulnerabilities that privilege-only scoring models
cannot detect.

AgentSentry is available as open-source software at
\url{https://github.com/Abhiram-ops/agent-sentry}.

% ──────────────────────────────────────────────────────────────────────────────
\begin{thebibliography}{00}

\bibitem{crowdstrike2024}
CrowdStrike, ``2024 CrowdStrike Global Threat Report,'' CrowdStrike, Tech. Rep., 2024.

\bibitem{langchain2023}
H. Chase, ``LangChain,'' 2023. [Online]. Available: \url{https://github.com/langchain-ai/langchain}

\bibitem{crewai2024}
J. Moura, ``CrewAI: Framework for Orchestrating Role-Playing Autonomous AI Agents,'' 2024.
[Online]. Available: \url{https://github.com/joaomdmoura/crewAI}

\bibitem{autogen2023}
Q. Wu et al., ``AutoGen: Enabling Next-Gen LLM Applications via Multi-Agent Conversation,''
\textit{arXiv preprint arXiv:2308.08155}, 2023.

\bibitem{mitre2024}
MITRE Corporation, ``MITRE ATT\&CK Enterprise Matrix,'' 2024.
[Online]. Available: \url{https://attack.mitre.org}

\bibitem{owasp_llm2023}
OWASP, ``OWASP Top 10 for Large Language Model Applications,'' 2023.
[Online]. Available: \url{https://owasp.org/www-project-top-10-for-large-language-model-applications/}

\bibitem{perez2022}
F. Perez and I. Ribeiro, ``Ignore Previous Prompt: Attack Techniques for Language Models,''
\textit{arXiv preprint arXiv:2211.09527}, 2022.

\bibitem{greshake2023}
K. Greshake et al., ``Not What You've Signed Up For: Compromising Real-World LLM-Integrated
Applications with Indirect Prompt Injection,'' \textit{arXiv preprint arXiv:2302.12173}, 2023.

\bibitem{yi2023}
J. Yi et al., ``Benchmarking and Defending Against Indirect Prompt Injection Attacks on
Large Language Models,'' \textit{arXiv preprint arXiv:2312.14197}, 2023.

\bibitem{sheyner2002}
O. Sheyner et al., ``Automated Generation and Analysis of Attack Graphs,''
in \textit{Proc. IEEE Symposium on Security and Privacy}, 2002, pp. 273--284.

\bibitem{saltzer1975}
J. H. Saltzer and M. D. Schroeder, ``The Protection of Information in Computer Systems,''
\textit{Proc. IEEE}, vol. 63, no. 9, pp. 1278--1308, 1975.

\bibitem{aws_iam2024}
Amazon Web Services, ``AWS Identity and Access Management User Guide,'' 2024.
[Online]. Available: \url{https://docs.aws.amazon.com/IAM/latest/UserGuide/}

\bibitem{azure_identity2024}
Microsoft, ``Managed Identities for Azure Resources,'' 2024.
[Online]. Available: \url{https://learn.microsoft.com/en-us/azure/active-directory/managed-identities-azure-resources/}

\bibitem{k8s_sa2024}
Kubernetes, ``Configure Service Accounts for Pods,'' 2024.
[Online]. Available: \url{https://kubernetes.io/docs/tasks/configure-pod-container/configure-service-account/}

\bibitem{conjur2022}
CyberArk, ``Conjur Open Source,'' 2022.
[Online]. Available: \url{https://www.conjur.org}

\bibitem{workload_identity2023}
Google Cloud, ``Workload Identity Federation,'' 2023.
[Online]. Available: \url{https://cloud.google.com/iam/docs/workload-identity-federation}

\bibitem{access_analyzer}
Amazon Web Services, ``AWS IAM Access Analyzer,'' 2024.
[Online]. Available: \url{https://aws.amazon.com/iam/access-analyzer/}

\bibitem{orca2023}
Orca Security, ``Cloud Attack Path Analysis,'' 2023.
[Online]. Available: \url{https://orca.security}

\bibitem{wiz2023}
Wiz, ``Cloud Security Graph,'' 2023.
[Online]. Available: \url{https://www.wiz.io}

\end{thebibliography}

\end{document}
